\chapter{内外挤迫下的社会政治困境}


物质和经济资源是第五次反“围剿”中中共遭遇的最大难题。与此同时，中共向来可资依赖的良好的政治资源也呈逐步削弱之势。在苏区内部，随着压力的加大，政治领导人年轻和缺乏经验的弱点日渐暴露，初期粗放发展阶段可能被隐藏和忽略的问题集中显现，在肃反、扩红、政权建设、群众支持等问题上，都出现不容乐观的态势。前方战事的不利，又像滚雪球一样，不断放大着这一切，而这些反过来又进一步加剧着军事的紧张。恶性循环的结果，使军事和政治的双重困境，降临到承受重重重压的中央苏区头上。

\section{肃反问题}


肃反是苏区时期一个沉痛的话题。1930年代前，中共的成长壮大，与苏俄的影响、帮助息息相关。这其中，是非纠缠、恩怨参杂。苏维埃发展过程中倍引争议的肃反运动，直接来源于苏俄的肃反理论及实践，\footnote{王首道回忆：“当时党中央派去各苏区的领导干部，传达贯彻苏联搁别乌（即保卫局）的肃反经验。”见王首道《回忆湘赣苏区》，《湘赣革命根据地》（下），第852页。}而且令人印象深刻的是，虽然处于不同的环境和发展阶段，各苏区肃反中出现的问题呈现相当的一致性，而苏区肃反的种种问题又和苏俄方面大同小异。


苏区艰难建设过程中，确实存在着各种各样的敌对势力，肃反有其逻辑上的合理性。但由于理论讹误、经验不足，加上战争环境的恶劣，各苏区在肃反过程中都存在过高估计敌对力量、采用逼供信恐怖手段、无视法律和事实等种种严重错误。有的地方甚至出现“宁肯杀错一百，不肯放过一个之谬论”。\footnote{《中共江西苏区省委四个月（一月至四月）工作总报告》，《江西革命历史文件汇集（1932年）》（一），第174页。}肃反的发展，常常都是向自己阵营的内部延伸，形成内部互相怀疑、自我损耗的局面，赣西南的反“AB”团、闽西的整肃社会民主党都是错误肃反的典型例证。赣西南因错打“AB”团激起红军兵变即“富田事变”；闽西肃社会民主党造成许多地区“党团内已抓了十分之七八”。\footnote{《少共闽粤赣省委一年来的工作总结》，《闽粤赣革命历史文件汇集（1932～1933）》，第16页。}盲目肃反的结果，瑞金全县“只有县委三四人，区委支部小组都没有了”。\footnote{《赣东特委步青给中央的报告（1931年8月8日）》，《中央苏区第三次反“围剿”》，第79页。}龙岩拥有七百多人的团组织“完全塌台”。\footnote{《少共闽粤赣省委一年来的工作总结》，《闽粤赣革命历史文件汇集（1932～1933）》，第14页。}湘鄂赣肃反甚至造成这样的情形：“区委向县委报告工作是隔一个山头望，不敢见面，怕杀掉了。隔着山头就喊，我那个地方发展了多少党员，搞了多少军队，最后总讲一声我是一个好人。”\footnote{傅秋涛：《关于湘鄂赣边区内战中期后期历史情形报告》，《湘鄂赣边区史料》，第54页。}肃反的错误既包括苏维埃内部的错误整肃，也包括对被认为是对立面的无原则打击，如湘赣省的酃县“把十六岁以上卅岁以下豪绅家属的壮丁无论男女都杀掉了”，\footnote{《中共湘赣苏区省委报告》，《湘赣革命根据地》（上），第271页。}这种肉体消灭政策虽在苏俄乃至许多革命运动中屡见不鲜，但和中共的基本方针仍然背道而驰。


中共中央领导人陆续进入中央苏区后，客观看，经过早期肃反的恐怖后，对肃反中的问题有所认识，并采取了一些纠正措施，红军和苏维埃政权内部大规模的肃反恐怖得到遏制。1931年12月13日，中华苏维埃共和国中央执行委员会非常会议发布关于《处理反革命案件和建立司法机关的暂行程序》的第六号训令，规定：

\begin{quote}

一切反革命的案件，都归国家政治保卫局去侦察、逮捕和预审，国家政治保卫局预审之后以原告人资格，向国家司法机关（法院或裁判部）提出公诉，由国家司法机关审讯和判决。

一切反革命案件审讯（除国家政治保卫局外）和审决（从宣告无罪到宣告死刑）之权，都属于国家司法机关。县一级司法机关，无判决死刑之权……中央区及附近的省司法机关，作死刑判决后，被告人在十四天内得向中央司法机关提出上诉。在审讯方法上，为彻底肃清反革命组织，及正确的判决反革命案件，必须坚决废除肉刑。\footnote{《中华苏维埃共和国中央执行委员会训令第六号》，《中央革命根据地史料选编》（下），第657～658页。}
\end{quote}


训令强调肃反要依靠专职机关和司法程序，把死刑判决权上收，防止基层乱抓乱杀。随后，苏区领导机关对前期肃反中暴露的严重问题陆续作出组织处理，纠正部分明显的错案。1933年，又专门发出训令，强调：“绝对废止肉刑，区一级裁判部不经上级裁判部的特许，绝对不许随便杀人。”\footnote{司法人民委员部：《对裁判机关工作的指示（1933年6月）》，石叟档案008·548/3449/0759，中国社会科学院近代史研究所藏。}中共中央采取的一系列制止盲目肃反的政策，缓和了富田事变前后中央苏区形成的恐怖肃反局面。


但是，过火的肃反思路形成并非一朝一夕，战争环境下这种思路更有其生长的土壤。特别是第五次反“围剿”开始后，随着战局向着不利于红军的方向发展，中共中央对肃反的判断再次严峻。在苏区遭受包围的战争形势下，中共中央对苏区内部的政治力量对比仍然缺乏足够信心，而苏维埃运动本身出现的一些问题造成民众的对立情绪更加剧了其形势紧张的判断。以此，中共中央过于悲观地判断苏区内部阶级关系，夸大苏区内部的敌对势力。中共中央领导人公开表示：“地主阶级在新区边区，特别在白军与刀团匪骚扰的区域，我们不但要在经济上消灭他们，而且要尽量在肉体上消灭他们。”\footnote{张闻天：《闽赣党目前的中心任务》，《斗争》第71期，1934年9月7日。}这等于肯定了不以事实而单以阶级划分进行肉体消灭的恐怖行动。在此观念指导下，肃反中的问题依然存在并发展，打击面涉及苏区社会的各个阶层。1933年6月苏区展开查田运动后，地主、富农出身人口在中央苏区普遍上升到总人口的10%以上，加上被作为打击对象的商人、资本家及其代理人、宗教人士、阶级异己分子、反革命、刀团匪，所谓敌对力量的人员空前增加，形成处处皆敌的局面；同时，中共中央对苏维埃政权内部也缺乏必要的信任，常常把查田运动、扩红运动遇到的问题乃至群众逃跑事件归结为苏维埃内部暗藏的“地主资产阶级的分子和他们的走狗”的破坏，甚至判断：“在我们党与苏维埃机关内埋伏着的‘坏蛋’不在少数。”\footnote{洛甫：《把革命的警觉性更加提高起来》，《斗争》第41期，1934年1月5日；《对于我们的阶级敌人，只有仇恨，没有宽恕》，《红色中华》第193期，1934年5月25日。}要求在政权内部进行广泛的检举、清查运动。


1934年初，第五次反“围剿”激战正酣时，中共中央决定展开进一步肃反的检举运动。3月21日，苏维埃中央政府在动员开展检举运动的大会上，推举产生9人组成的检举委员会，领导整个中央苏区的检举运动。4月8日，临时中央政府中字第五号《命令》明确规定废止1931年底至1932年初为纠正肃反错误而制定的将死刑判决权上收、严格司法程序等条例，表明中共中央准备采取严厉的镇压措施。4月19日，《红色中华》发布中央工农检查委员会第二号训令，指出目前敌人五次“围剿”的决战已到了最紧张最尖锐的决定最后胜负的阶段，为保证战争的胜利，必须检举“苏维埃机关内的消极怠工的分子，贪污腐化、浪费的分子，脱离群众离开群众利益和工作上的官僚分子，退却逃跑、动摇不坚定的分子，包庇地土富农与妥协的分子，违反法令与破坏纪律的分子，特别是检举暗藏在苏维埃机关内的阶级异己分子和反革命。”由此，检举运动在中央苏区开展起来。


为推动运动的深入发展，1934年5月12日中共中央发出《给各级党部党团和动员机关的信》，提出：“敌人进攻更加紧张，反革命活动也更加厉害。”“因此，我们必须更加提高阶级警觉性来对付各种各式的、埋伏的或公开的反革命分子。对于各个机关和赤少队，必须继续进行检举，不仅检举其成分，特别要注意从政治上检举。检举机关、裁判部、保卫局的系统，必须活跃起来。但突击队须更负责的领导肃反的工作，不要机械的等待保卫局裁判机关与检举机关。”\footnote{《给各级党部党团和动员机关的信》，《中共中央文件选集》第10卷，第269页。}同时，中共中央领导人连续发表文章，号召展开严厉肃反，强调：

\begin{quote}

不认识苏维埃法庭是阶级斗争的工具，是压迫敌对阶级的武器，而表现出单纯的法律观，机械的去应用法律。不知道法律是随着革命的需要而发展，有利于革命的就是法律。凡是有利于革命的可以随时变通法律的手续，不应因法律的手续而妨碍革命的利益……许多裁判机关，侧重于法律手续，机械的去应用法律、对镇压反革命的重要工作却放松了，这是裁判机关在工作上的极大缺点，而且是严重错误……清理档案，凡有反革命事实的豪绅地主富农等阶级异己分子，经公审后，立即执行枪决。\footnote{梁柏台：《裁判机关的主要工作方向——镇压反革命》，《红色中华》第156期，1934年3月1日。}
\end{quote}


《红色中华》则发表社论要求：“在战区边区，我们对于任何反革命的活动，必须立刻采取最迅速的处置，凡属进行反革命活动的豪绅地主、富农、商人、资本家、老板、流氓，必须立刻捉起。除个别最重要的分子须严究同党外，其余无须详审，无须解县，一概就地枪决。就是他们中间有反革命嫌疑的分子，也应即刻捉走，重的当地枪决，轻的押解后方监禁。”社论以异常严厉的口气强调：“一切对于反革命的宽容与放纵，一切‘讲究手续’与‘法律观念’，一切犹豫不决与迟缓，在目前同阶级敌人决死战的时候，客观上都是反革命的助手与帮凶。”\footnote{《对于我们的阶级敌人，只有仇恨，没有宽恕》（社论），《红色中华》第193期，1934年5月22日。}


检举运动铺开后，中央苏区清洗了一大批干部，万泰“检举出了县苏军事部长、收发员、窑下区苏主席是AB团，破获了AB团的县团部，并且还检举出了县内务部的科员包庇地主偷公章，少共县委总务科长偷路条并且贪污”。\footnote{《在加紧整理中的万泰工作》，《红色中华》第191期，1934年5月21日。}运动后期，《斗争》曾发表文章总结运动成果：

\begin{quote}

中央政府各部共洗刷了六十四人，其中有九个贪污的，十五个破坏苏维埃法令和政府威信的，四十个消极怠工自由回家的，江西乐安县一级及善和增田两区乡共洗刷了六十二人，万太县区各机关中洗刷了六十人，胜利县一级二十三人，石城县一级洗刷了二十八人。区一级二十五人、乡一级四人……粤赣于都自进行检举以来，在县一级洗刷了三十八人，在区乡共五十五人，合作社洗刷了三十一人；西江在县一级洗刷了二十人，区乡共六十人；会昌检举才开始，已在县一级洗刷了九人。\footnote{董必武：《把检举运动更广大的开展起来》，《斗争》第61期，1934年5月17日。}
\end{quote}


总结中谈到的数据主要还是针对领导干部的，社会上的清洗尚不在其列，即使如此，已可看到其牵涉之大。


检举运动中，扩大化的事例处处可见。1934年5月，西江一县在“下半月短短的半个月中，即捕获了几百名反革命分子，只判处死刑的即有二三百名（城市区在红五月中共杀了三十二名反革命，破获了AB团、暗杀团、铲共团、社民党、保安会的组织，共捉了四个暗杀团长，两个AB团长，数十名连长、排长、宣传队长等）”。\footnote{《西江县——红五月扩红突击中的第二名》，《红色中华》第199期，1934年6月7日。}西江是由瑞金、于都、会昌划属的小县，人口仅数万人，半月内即出现如此之多形形色色的“反革命”和反革命组织，这本身已极不正常，而这种做法还作为正面典型受到鼓励。瑞金“一个七十岁的贫农，在闲谈中说到白军到了清流归化，却判决了死刑”。\footnote{司法人民委员部：《对裁判机关工作的指示（1933年6月）》，石叟档案008·548/3449/0759，中国社会科学院近代史研究所藏。}闽赣省裁判部的钟光来甚至“把裁判部犯人大批的不分轻重的乱杀一顿”，在由樟村退往石城途中，更是“沿途格杀群众”。\footnote{《闽赣省枪毙反革命首领两只》，《红色中华》第180期，1934年4月26日。}恐怖的气氛，“使一部分被欺骗群众首先是中农群众登山逃跑，或为地主富农所利用来反对苏维埃政权”，催生出“因为肃反工作的加紧而逃跑到高山上去了的反革命武装”。\footnote{张闻天：《是坚决的镇压反革命还是在反革命前面的狂乱？》，《红色中华》第208期，1934年6月28日；《西江县——红五月扩红突击中的第二名》，《红色中华》第199期，1934年6月7日。}肃反扩大化的问题，正如张闻天所反省的：“一些地方，赤色恐怖变成了乱捉乱杀，‘阶级路线’与‘群众路线’也不讲了。在一些同志中间正在形成‘杀错一两个不要紧’或者‘杀得愈多愈好’的理论……对于肃反中的恐怖主义，那更是没有人敢讲话，因为恐怕批评这种倾向时，人家就会把他当作反革命的同道者看待，而性命难保。”\footnote{张闻天：《反对小资产阶级的极左主义》，《斗争》第67期，1934年7月10日。}


现存的一些案例真实反映出当时中共中央在肃反问题上的偏差。1933年，江西宜黄判决的几个反革命案例分别是，卢章秀：“（一）该犯任黄陂区委少队部的时候，五月间，往中央总队部开会，自愿报名加入少共国际师，结果回来向县队部请假回家，将近两个多月，不但不去少共国际师当兵，而且坚决不来工作。（二）不经过介绍，自由行动去保卫局工作。”曾彦贵：“该犯于今年三月间到保卫局当看守兵，因看守不注意，走了二个靖卫团，就坐了十余天禁闭。（二）后放他出来，后在保卫局当挑夫，一贯的消极怠工。（三）叫该犯去买米，故意贪污大洋一元。”罗宦泉：“（一）该犯联名孟章，坚决担保反动富农（陈国芳妻），代表土豪说没有钱，结果把（陈国芳妻）放出来，随即跑下宜黄城去了。（二）我红军独立营捉获一个反革命，该犯又亲自来担保。”\footnote{《江西省宜黄县裁判部法庭判决书，第5号》，石叟档案008·548/3449/0759，中国社会科学院近代史研究所藏。}从文件看，上述几个案犯都没有确切的反革命罪的证据，列举的问题，不足以构成事实上的犯罪，将其定为反革命犯，反映出法律观念的淡漠。正如于都报告的，在打击反革命时，“捉了一些造谣的，都没有多大事实”。\footnote{《中共于都县委七、八、九三个月工作总结》，《江西革命历史文件汇集（1932年）》（二），第257页。}


由于政治凌驾法律，定罪随意性强，冤假错案发生几率很高。于都的一个案例很具代表意义：新陂中塅乡丁福生“用恐吓的手段，对付红军家属丁昌早要量米给他吃，因没有满足他的敲索，便弄国民党证，到政府妄报丁昌早为国民党员，当时政府也不加以调查和侦察，听信所言的话，便把丁昌早家产没收，弄得丁昌早全家苦不堪言”。也许是丁福生从这种政治陷害中尝到了轻易得手的甜头，随后他“又与丁昌钊因为耕牛的嫌隙，妄报丁昌钊为地主，其实丁昌钊经调查确系中农”。丁福生的行径固然恶劣，但仍属于挟嫌报复的诬告行为，并不一定具有破坏革命的主观企图，但是当丁的诬告行径暴露后，却被定性为：

\begin{quote}

丁福生的这种企图，是动摇红色战士，故意破坏苏维埃联合中农的策略，有意帮助反革命的实际行动，是苏区里面的内奸。此事幸经发觉，经于都县裁判部判处该犯以监禁半年，送交省裁判部审查，认为该犯妄报红军家属为国民党员，是动摇红军军心，妄报中农为地主，是破坏苏维埃对中农的政策，很明显的是反革命作用，应判处死刑，已批复该县遵照批示执行。\footnote{《挟嫌妄报的反坐》，江西省苏维埃政府裁判部《司法汇刊》第1期，1933年6月16日。}
\end{quote}


这样的判决，以判断代替事实，处理显然失之草率和严苛。


事实上，由于反革命罪是一个很宽泛的概念，当时许多触及法律者被加以反革命罪论处。于都岭背区的谢锦波、谢正月生等十余人“受查田运动浪潮的冲激，便主谋逃往白区”，他们联系了一批人，其中有人报告了苏维埃政权，结果这些人“行经梓山地方过，就被区府预先埋伏在那里的革命武装，通统捉着，一个也没逃走……在群众的呼声喊杀之下，对谢锦波、谢正月生和豪绅地主富农总共十一名都判处以死刑，并就地执行枪决了。”\footnote{《于都岭背区开始进行查田运动中》，《司法汇刊》第2期，1933年7月9日。}惩处这一案件中的首犯自无问题，但重判理由、范围及就地执行的方式其实都缘于反革命这一定罪，而一旦被视为反革命，在当时的环境下，结局就几乎决定了。张闻天曾明确表示：“在某种条件之下，从法律上说来某个反革命分子枪决的法律根据还没有找到，但是在群众的热烈要求枪决的条件之下，我们把他拿来枪决，以满足群众的要求，发动群众的斗争，还是为我们所容许的。”甚至强调：“必须经常使用群众的暴力去与反革命作斗争。一切法律观念是极端有害的，甚至是自杀的。”\footnote{张闻天：《无情地去对付我们的阶级敌人》、《闽赣党目前的中心任务》，《斗争》第49期，1934年3月2日；第71期，1934年9月7日。}


在加紧肃反的思路指导下，执法机关出于慎重对一些案件进行的调查取证工作被指为“机械的法律观念”，“客观上都是反革命的助手与帮凶”。\footnote{《对于我们的阶级敌人，只有仇恨，没有宽恕》（社论），《红色中华》第193期，1934年5月25日。}公略县裁判部长对案件处理较为慎重，重大案件强调送上级批准和材料充分，即被指责为“浓厚的机械的法律观念”，\footnote{《在整理裁判部工作中中央司法部洗刷动摇妥协分子》，《红色中华》第174期，1934年4月12日。}作为动摇妥协分子典型受到严厉批评。苏区中央领导人公开表示：“不必需要多少法律的知识，只要有坚定的阶级立场，他就可以正确的给犯罪者以应得的处罚。”要求：“以后的案件，应随到随审，非有特别事故，自受到案件之日起，不得超过三天，就要解决。”仅仅是有“反革命嫌疑的分子，也应即刻捉起，重的当地枪决，轻的押解后方监禁”。\footnote{洛甫：《无情的去对付我们的阶级敌人》，《斗争》第49期，1934年3月2日；梁柏台：《裁判机关的主要工作方向——镇压反革命》，《红色中华》第156期，1934年3月1日；张闻天：《对于我们的阶级敌人，只有仇恨，没有宽恕》，《红色中华》第193期，1934年5月25日。}领导人这样的认识加上中央苏区本就薄弱的法律背景，使法律运用难以健全：“有些地方（会昌、石城），审判案件不是在法庭上公开审判，而是在裁判部长的房间里，甚有将处死刑的案件，未经过法庭审判，在房间里写个判决书，送上级去批准执行，群众不知道究竟为什么事情杀人。”\footnote{司法人民委员部：《对裁判机关工作的指示（1933年6月）》，石叟档案008·548/3449/0759，中国社会科学院近代史研究所藏。}同时，逼供成为审讯的重要方式，虽然苏区有关法令明令“废止杀头破肚及肉刑等刑罚”，但又规定“为取得犯人实供，如敌探等有时得用肉刑讯问”。\footnote{《裁判条例》，《福建革命历史文件汇集》甲14册，第83页。}1933年初，中共湘鄂赣省委也提出：“对审讯犯人固然要纠正‘左’倾的单凭刑询的错误，但是认为刑询便是‘左’倾亦是另一个极端的右倾错误。”\footnote{《中共湘鄂赣省委第二次执委会决议案（1933年1月3日）》，《湘鄂赣革命根据地文献资料》第2辑，第15页。}这实际是在为刑讯手段开方便之门。


肃反中对打击面的任意扩大，一定程度上影响到中共与民众间原有的信任关系，当时，在苏区一些地区“看到反动标语，似乎并不算一回事”。\footnote{张闻天：《于都事件的教训》，《斗争》第53期，1934年3月31日。}边区有些地方国民党组织的民团武装“敢长驱直入的到四边围绕有赤区的区政府捉人、缴枪”，而“群众对此事好象没多大关系一样”。\footnote{童小鹏：《军中日记》，1933年8月21日，第41页。}更严重的是：“一部分被欺骗群众首先是中农群众登山逃跑，或为地主富农所利用来反对苏维埃政权。”\footnote{张闻天：《是坚决的镇压反革命还是在反革命前面的狂乱？》，《红色中华》第208期，1934年6月28日。}对地主肉体消灭的做法甚至到1949年中共横扫江南时仍使一些人心有余悸：“江西许多地区的地主、富农，因受过去十年内战时期土改的偏向影响，误认为划成地富不但是多要粮，而更重要的是要命的问题。”\footnote{《江西省人民政府财政厅6个月（从初解放到年底）工作综合报告》，程懋玲编《江西财政报告集》，江西人民出版社，1991，第4页。}其负面影响不可谓不深。

\section{工作作风问题的滋生}


中央苏区和全国各苏区一样，是在实际的革命运动中成长壮大的。在这一过程中，如毛泽东本人的工作作风所体现的，注重实际、强调调查研究、不务虚文是其明显的特点。中共中央领导人陆续进入中央苏区后，着力进行正规化建设，在组织建设和干部管理上采取一系列措施，使苏区政治再上一个台阶。不过，作为一个崭新的革命政权，在形成新的政权架构、行政理念和工作作风时，中共仍然面临着相当多的考验。苏区在建立强有力的国家机器的同时，缔造出具有支配权力的强势管理者，而苏区的现实状况又完全符合列宁所说的官僚主义的经济温床：“小生产者的分散性和涣散性，他们的贫困、不开化，交通的闭塞，文盲现象的存在，缺乏农工业之间的流转，缺乏两者之间的联系和协作。”\footnote{列宁：《论粮食税》，《列宁全集》第41卷，人民出版社，1984，第218页。}在这样的环境下，正如列宁所说十月革命后一年多时间，“官僚主义就在苏维埃制度内部部分地复活起来”\footnote{列宁：《论粮食税》，《列宁全集》第41卷，第217页。}一样，在苏区它也很难避免地会成为吞噬苏维埃制度活力的怪兽。而中共中央具体指导中，在加强正规化同时对所谓“经验主义”的批评，又给了形式主义、官僚主义更多的养分。潘汉年曾形象描绘苏区基层党支部开会走过场的情形：

\begin{quote}

参加支部会议的方式，多半是学得县委，做一个又长又臭的报告，弄得支部同志都不肯或不敢发意见，一连催几声：“同志们！话！话！”假如还没有人说，支部书记或参加的区委，便提高喉咙问一句：“大家听懂没有？”下面齐声答应：“听懂了！”上面再问：“冒马格意见？”“冒！”或者再来一问：“同意唔同意？”又是齐声响亮的回答“同意！”\footnote{潘汉年：《这样的工作作风好不好》，《斗争》第63期，1934年6月9日。}
\end{quote}


在苏维埃胜利发展时期，上述问题还不是那么引人注目。第五次反“围剿”开始后，由于军事不利，苏区社会经济环境恶化，不利的形势加剧了苏区政治中的一些问题，而且使其影响更加放大，最明显的反映就是为支持前方战事而在后方的强迫命令。


1933年后，为应对国民党军的封锁，应付战争需要，中共中央不得不大规模发行公债，向民众大量借谷，使地方政权面临很大压力。有关研究显示，1934年中央苏区农民负担平均达到其收入的15.7%，有的甚至达到30%以上，“不能不承认负担是重的”。\footnote{《中国农民负担史（第3卷）中国新民主主义革命时期革命根据地的农民负担（1927～1949）》，中国财经出版社，1990，第105页。}沉重的负担使已经饱经战争摧残的苏区民众难以承受。1933年8月苏维埃中央决定发行300万元经济建设公债，到11月15日江西实际收得现款仅42万余元，不到江西全省领去的公债数的20%，其中杨殷、南丰、太雷、长胜、崇仁等县未收到一分钱。\footnote{《中共江西省委关于全省推销经济建设公债的初步总结》，《江西革命历史文件汇集（1933～1934年）》，第305页。}在苏区行政效率相当高的背景下，这一状况的出现主要是由于民众财力已告空乏。但当时的苏区领导者对此缺乏足够认识，为维持军队和政权的运转，仍不断推出一些过高的指标与要求，并以机会主义、动摇倾向指责没有完成高指标者，使各地为完成指标而强迫命令成风。万泰县冠朝区十个乡，公债销售中“摊派的七个乡”；“会寻安有几个乡扩大红军成为群众的恐怖，听到工作人员下乡，就纷纷上山或躲避不见”。\footnote{陈寿昌：《万泰工作的转变在哪里？》，潘汉年：《工人师少共国际师的动员总结》，《斗争》第39期，1933年12月19日；第24期，8月29日。}当时，这些问题绝非个别，正如苏维埃中央反省的：“坐禁闭，罚苦工，差不多是这些工作人员对付群众的唯一办法”。\footnote{《人民委员会为万太群众逃跑问题给万太县苏主席团的指示信（1934年4月3日）》，《红色中华》第137期，1934年4月10日。}


客观而言，在战争形势下，加强对资源的吸取，有其可理解的不得已之处，但是，忽视民众必要的利益要求，一味要求民众作出牺牲，则不免有竭泽而渔之嫌。而将反对、抵制这种做法的干部或仅仅是提出意见者都视作异己分子加以批判，无异是在自毁墙脚。远在莫斯科的王明通过中共赴莫斯科代表了解到苏区实况后，对此提出批评，指出中共中央“党内两条路线斗争”名义下存在的“不可忽视的严重的弱点”：

\begin{quote}

A.对于缺点和错误的过分和夸大的批评，时常将个别的错误和弱点都解释成为路线的错误……决没有领导机关的路线正确，而一切被领导的机关的路线都不正确的道理，此种过分和夸大的批评，既不合乎实际，结果自不免发生不好的影响，一方面不能真正推动工作，另一方面使地方党部的工作人员发生害怕困难，对困难投降的情绪，而且甚至使一部分幼稚的同志发生跳不出机会主义的泥坑的烦闷心理，以致有的发生对党和革命抱悲观失望的态度。B.对于党内斗争的方法有时不策略，比如在中央苏区反对罗明路线时，有个别同志在文章上，客观上将各种的错误，都说成罗明路线的错误，甚至于把那种在政治上和个人关系上与罗明路线都不必要联系在一起的错误，都解释成罗明路线者。这样在客观上不是使罗明孤立，而恰恰增加了斗争中可以避免的纠纷和困难。\footnote{王明、康生：《给中共中央政治局的信（1934年4月20日）》。}
\end{quote}


王明的批评可谓中肯，动辄以“路线错误”责人，是政治正确逻辑下官僚主义的另类生存，如此盲目上纲的结果只能造成人们谨小慎微，不敢越雷池一步。李维汉当时曾撰文批评道：“某某省工会和县委的负责同志，误解思想斗争是禁止别的同志对于工作上实际问题发表不同的意见。干部间在共同一致为要完成中央给予突击任务的时候，对于实际办法与领导同志有不同意见时，一概拒绝，甚至批评他们反党。结果许多干部见了某某省工会和县委的负责同志，不敢说话，只是‘唯唯而退’。”\footnote{罗迈：《把突击运动期间党内斗争上表现出来的缺点与错误纠正过来》，《斗争》第51期，1934年3月17日。}人们后来回忆，由于怕犯错误，被打成反革命，“那个时候做工作只能在一定的范围内，出了范围就不行，怕犯错误。这是在那个时期的现象”。\footnote{傅秋涛：《关于湘鄂赣边区内战中期后期历史情形报告》，《湘鄂赣边区史料》，第54页。}


强迫命令、官僚主义现象的盛行，除苏区政治经济背景及上层指导的错误外，基层组织软化也起了催化作用。由于大批优秀基层干部输送前方及中共组织发展中的“唯成分论”观念，苏区地方干部素质下降严重，1933年江西县一级419名干部中，能读会写的154名，只占37%。\footnote{《党的组织状况》，《中央革命根据地史料选编》（上），第687页。}虽然各级政权采取了一些措施，提高基层干部素质，但多数是形式重于内容。支部会议往往开成这样的结果：“支部每次开会，是非常机械的，对支部实际问题与群众斗争，很少讨论，只是出席人及支书作一个报告，支部同志没有意见就完了。在讨论中只是由速记来长篇抄写上级决议或由速记作几条符咒式的简单的决议，到会的同志被剥夺发言权似的纷纷‘开小组会’、‘学时文’，甚至打瞌睡去了，主席一声‘通过决议’，会场陡然壮严肃静，一声‘同意’，万事皆休，从报告到结论都是出席人或速记的声音。”\footnote{《中共湘赣苏区省委报告》，《湘赣革命根据地》（上），第602页。}在这些因素影响下，各地苏维埃政权组织、精神和活动能力堪忧，石城县反映：“县政府的一些干部发生不好的现象，如派去工作开小差回家的，派去扩大红军突击工作，带红军家属老婆到乡苏睡觉……一般干部特别是区一级的干部出发工作，就开小差回家了，在机关里吃饭时就来，吃了饭就回家去，找人来工作找人不到。”\footnote{《梁广给全总执行局的报告》，转见刘少奇《反对扩大红军突击运动中的机会主义的动摇》，《斗争》第41期，1934年1月5日。}胜利马安石公布乡“支部书记与乡苏主席都是官僚主义的标本，什么工作只说下了条子通知了，自己一点事不理，半步门不出。”\footnote{《中共江西省委红五月工作总报告（1933年6月17日）》，《江西革命历史文件汇集（1933～1934年）》，第149页。}还有一些干部则利用权力“在乡村中作威作福，无所不为”\footnote{《人民委员会为万太群众逃跑问题给万太县苏主席团的指示信》，《红色中华》第137期，1934年4月10日。}。


令人惊讶的是，中央苏区后期基层干部大吃大喝现象已不罕见。江西省苏1932年5月披露：“一个乡政府以前经常可有三四桌人吃饭，区政府可有七八桌人吃饭，每月的客饭一个乡政府可开一二百元。”\footnote{《江西苏区中共省委工作总结报告（1932年5月）》，《中央革命根据地史料选编》（上），第450页。}闽西永定县“有些政府好象变成了一些客栈、饭店，无论什么人都可以在政府食饭”，宁化县“县主席不知常驻工作人员有多少”，“只见人吃饭，不见人工作”。到饭馆公款吃喝也有案可稽，《红色中华》揭露：“靠近瑞金县苏的一家菜馆，据说是专供瑞金县苏工作人员食的，有一次购买十六斤甲鱼，还说不够卖”。\footnote{《开展广泛的反贪污斗争》，《红色中华》第134期，1933年12月11日。}大吃大喝的背后通常都伴随着贪污浪费，江西省苏的报告指出：“各级政府浪费的情形实可惊人，一乡每月可用至数百元，一区一用数千，一县甚至用万元以上；贪污腐化更是普遍，各级政府的工作人员随便可以乱用隐报存款、吞没公款，对所没收来的东西（如金器物品等）随便据为己有，实等于分赃形式。”\footnote{《江西苏区中共省委工作总结报告（1932年5月）》，《中央革命根据地史料选编》（上），第450页。}


一些地方干部特殊化现象也在滋生，即便像优待红属这样严肃的政治活动也有人借机营私：“宁化各级苏维埃压迫群众替自己工作人员作过分的劳动，而红军家属反得不到优待。”\footnote{严仲：《三个月扩大红军的总结与教训》，《中央革命根据地史料选编》（中），第664页。}宁都某区政府“对自愿的当红军的反不替他耕田，并且说谁要你去当红军，而对强迫命令去当红军的才替他耕田”。\footnote{《中共江西苏区省委四个月工作总结报告（1932年5月）》，《江西革命历史文件汇集（1932年）》（一），第138页。}之所以如此，关键在于自愿当红军的不属指标范围。博生县苏工作人员的优红工作严重流于形式：

\begin{quote}

某天到城市流南区流南乡去实行“礼拜六”，大概是九点钟才从县苏起程，到流南乡已经差不多十点钟，找了乡政府和代表带到红军家属的地方已经十点半钟过了，在田野中还没做到两点钟的工夫，县劳动部有一个同志发了“摆子”，如是找着那村的乡代表，要他立即找担架将病人抬进城去。那位乡代表的回答是很实际的：“你们这样优待红军家属太形式了，没有做到两个钟头的工夫就要找担架，你们不来还免得我们麻烦。”\footnote{《在苏维埃系统内开展反官僚主义的群众斗争（1933年11月20日）》，中共江西省委：《省委通讯》第47期，1933年11月20日。}
\end{quote}


粗暴的工作作风和形式主义、官僚主义倾向，严重伤害了群众的感情。黎川县樟村“赤卫军连长要群众帮助红军运输，用命令强迫，结果引起了群众的反抗”。\footnote{童小鹏：《军中日记》，1933年10月20日，第51页。}于都“寨下区发生反革命收缴保卫局的枪支，捉区苏主席，苍前区发现反革命的标语，小溪罗江等区反革命分子能够回来，公开宣传要群众把已经买了的公债退还政府，强迫乡苏代表收回公债，反对推销公债”。\footnote{盛荣：《于都问题》，《青年实话》第3卷第15期，1934年3月。}时任粤赣军区司令员的何长工回忆，粤赣省的罗田甚至发生红军“打仗时刚过了一座桥，群众把桥给破坏”\footnote{《何长工回忆录》，解放军出版社，1987，第318页。}的恶性事件。这种情况虽不是普遍现象，但群众正面支持的减弱却是不争的事实，安远、寻乌有干部反映：“群众情绪更加消沉，红也不管，白也不管。群众受反动派压迫，一点也不会反抗而离开我们。”\footnote{《粤赣省委对寻乌安远工作的决定》，转见罗迈《在粤赣省第一次党代表会议的前面》，《斗争》第34期，1933年11月12日。}群众支持度的下降，对“只有动员群众才能进行战争，只有依靠群众才能进行战争”\footnote{毛泽东：《关心群众生活，注意工作方法》，《毛泽东选集》第1卷，第136页。}的红军展开反“围剿”作战，应该是不小的隐忧。

\section{群众逃跑}


1933年，中央苏区出现的群众成规模逃跑事件，是苏区发展过程中积累的一系列问题的集中显现。


群众逃跑在苏区初创时期也有出现，赣东北苏区初创时，“因经济发生恐慌，常有逃亡”。\footnote{《江西政治报告（1929年5月31日）》，《中央革命根据地史料选编》（中），第89页。}当时或是出于国民党长期负面宣传影响下对中共和红军的不理解、恐惧，或是对国民党军可能的报复行为的担心，或是苏区本身政策的失误，苏区或多或少有过部分群众逃跑。尤其是肃反的错误，曾造成群众短时期的恐慌和逃跑。1932年湘赣苏区报告，上犹县“营前一区逃跑了三千余群众到白区去”。\footnote{《红军三军团政治部关于崇犹两苏区路线和红军情况的报告（1932年5月15日）》，《湘赣革命根据地》（上），第290页。}不过总的来看，苏区稳定后，广泛涉及各种成分群众的大规模集中逃跑的事件较为少见。而且，还有群众往回流动的，毛泽东《兴国调查》中提到：“因为革命胜利，彭屋洞早先去泰和耕田的农民这时候回来了十二个人，因为泰和那边还没有革命，听到兴国革了命，有田分，都回来。”\footnote{《毛泽东农村调查文集》，第187页。}


中央苏区群众集中逃跑始于1933年下半年，首先从边区的万泰、于都、连城等地开始，并迅速蔓延。这一问题的出现，查田运动是直接导火索。何长工根据其在粤赣工作的经验谈道：“在地方工作中，实行了地主不分田，富农分坏田，限制中农发展的错误的土地政策，使部分群众发生动摇，根据地边沿地区出现了一个短时间的部分群众‘外逃’的现象。”\footnote{《何长工回忆录》，第313页。}早在运动大规模展开的初始阶段，一些地区就出现逃跑问题，乐安招携“在几天之内有几百群众随同富农地主跑到白区去”，寻乌、会昌等县在运动中查出的地主、富农纷纷“乘夜逃上山，实行土匪生活”。于都段屋、岭背、城市、寨下面等区“查出的地主土豪富农有三分之二乘着天雨水涨，星夜乘船顺水而逃，浸死颇多”。\footnote{《猛烈的开展查田运动》，中共江西省委《省委通讯》第20期，1933年7月9日；《豪绅地主的残余在查田运动中发抖》、《胜利于都地主富农企图逃跑》，《红色中华》第94期，1933年7月14日；第96期，1933年7月26日。}同时受到查田乱划成分影响，部分查田运动中尚未遭打击的普通民众也开始逃跑，胜利车头、河田、仙霞观等区“少数群众同下赣州城”。赣县长洛、大埠、白露、良口、大田等区“少数群众逃到白区”。\footnote{《胜利县查田运动的教训》，中共江西省委《省委通讯》第20期，1933年7月9日；《赣县查田运动胜利中的缺点》，《红色中华》第124期，1933年11月11日。}随着运动进一步展开，触犯中农、贫农的运动扩大化现象加剧，原来的贫农、中农大批被划为地主、富农，弄得人人自危，逃跑面迅速扩大，形成“成群结队整村整乡”\footnote{《人民委员会为万太群众逃跑问题给万太县苏主席团的指示信（1934年4月3日）》，《红色中华》第137期，1934年4月10日。此标题中的“万太”应为“万泰”。}逃跑的恶劣局面。于都“岭背区特派员乱打土豪，故意将中农当地主打，造成群众恐慌和逃跑”。\footnote{张闻天：《于都事件的教训》，《斗争》第53期，1934年3月31日。}万泰县窑下区郭埠乡不顾贫农团会议许多人反对，强行将一人划为富农，结果二三天内群众“就走了一大批”，由于该县普遍存在“工作人员乱打土豪，把贫农中农当做地主富农”的问题，群众逃跑十分严重，“塘上区有群众约六千人，逃跑的在二千人以上（一说二千三百人，县委报告是一千八百人），而且大部分是男子。”\footnote{陈寿昌：《万泰工作的转变在哪里？》，《斗争》第39期，1933年12月19日；《人民委员会为万太群众逃跑问题给万太县苏主席团的指示信（1934年4月3日）》，《红色中华》第137期，1934年4月10日；陈寿昌：《万泰工作的转变在哪里？》，《斗争》第39期，1933年12月19日。}资溪县由于“发展查田运动”，“十余天来，各区群众向白区逃跑现象日益发展，从一乡一区蔓延到很多区乡，从数十一批增加到几百以至成千人一路出去，从夜晚‘偷走’变而为明刀明枪的打出去，杀放哨的，甚至捆了政府主席秘书走。”\footnote{《关于资城事变问题省委对资溪县委的指示信》，石叟档案008·222/3745/0247。}


查田运动诱发群众逃跑，同时由于受到封锁，资源匮乏，当时的中共中央又缺乏应对危机的能力，高指标和政治威胁相结合，形式主义和强迫命令成风，进一步伤害了群众的感情，加剧群众逃跑现象。信康县牛岭区工作人员拿公债“挨屋挨户去摊发，使得群众不满意……少数落后的贫苦工农分子，跟着富农去反水”。\footnote{《显微镜下的官僚主义》，《红色中华》第131期，1933年12月2日。}1934年1月至3月，于都禾丰区就有600余群众逃跑。\footnote{《禾丰区破获反革命暗杀团》，《红色中华》第176期，1934年4月17日。}寻乌1933年初被国民党军短暂进攻，在恢复寻乌时，“全县反水群众有八千三百三十六人（以澄江、寻城为多），阶级异己分子约占百分之三十”，当时报告在谈到这一现象时客观承认：

\begin{quote}

大批的工农群众所以被敌人领导反水的原因，虽然是由于反革命的欺骗作用加强，但是过去寻乌群众阶级斗争没有深入，肃反工作不注意，工作方式命令欺骗强迫的结果，便利于敌人的欺骗恐吓，而使大批的工农群众脱离党及苏维埃的领导多是主要原因。\footnote{《中共寻乌县委一个半月动员工作总结报告（1933年3月4日）》，《江西革命历史文件汇集（1933～1934年）》，第32～33页。}
\end{quote}


由于战争的持续进行，苏区人力资源短缺问题凸显，巨大的兵员指标和苏区人力资源形成极大反差，为完成高额的扩红指标许多地区不得不采取强迫的办法，万泰县“冠朝区在一个会议上强迫一个反帝同盟主任报名，不报名就罚苦工当伙夫”。\footnote{陈寿昌：《万泰工作的转变在哪里？》，苏区中央局：《斗争》第39期，1933年12月19日。}由于此，逃跑成为民众躲避扩红的一种方法。于都“大部分模范赤少队逃跑上山，罗凹区十分之八队员逃跑上山，罗江区有300余人逃跑，梓山、新陂、段屋区亦发生大部分逃跑，有的集中一百人或二百人在山上，有的躲在亲朋家中”。该县新陂区密坑乡“精壮男子完全跑光了”。\footnote{《于都发生大批队员逃跑》，《青年实话》第111期，1934年9月20日；《于都在坚决执行党的指示中已经开始转变过来》，《红色中华》第238期，1934年9月26日。}


扩红运动中发生的种种极端行动造成干群间关系的严重紧张，瑞金白鹭乡“二十余名模范队员举行反动暴动，捆去区委人员三个，杀伤一个”。\footnote{《学习瑞金的经验与教训动员整营整连的模范赤少队武装上前线去》，《红色中华》第230期，1934年9月6日。}宁化县群众除以逃跑表示消极对抗外，还有公开反抗的行动，以致上级派出的突击队“吓得不敢出乡苏门口”。\footnote{《宁化落后的原因在哪里？》，《红色中华》第238期，1934年9月26日。}对此，中共中央曾作出检讨，坦承过失：“群众逃跑的主要原因是由于我们苏维埃政府领导上的错误。我们许多区乡苏维埃的工作人员，不论在推销公债，扩大红军或收集粮食方面都采取了严重的摊派与强迫命令的办法，任何宣传鼓励、解释说服的工作也没有。”\footnote{《人民委员会为万太群众逃跑问题给万太县苏主席团的指示信》，1934年4月3日，《红色中华》第137期，1934年4月10日。}这样的批评当然不是信口开河，竭泽而渔的做法实在也得不偿失，对此无论中共中央还是地方的实际执行者，应该都不会不清楚，只是面对兵员需求的巨大压力，他们实在也是左右为难。


不仅仅是普通民众，第五次反“围剿”期间，在工作遇到困难或国民党军进攻时，苏区干部开小差、逃跑现象也相当普遍：宜黄、西江、泉上负责肃反的裁判部长“公开的投降白匪”。甚至有领导群众逃跑反水的：“乐安、万太、广昌、代英、门岭、公略都发生过这种现象”。\footnote{梁柏台：《裁判机关的主要工作方向》，《红色中华》第156期，1934年3月1日；董必武：《把检举运动更广大的开展起来》，《斗争》第61期，1934年5月26日。}福建为培养地方武装领导干部开办的训练班，学员均经过各地选拔，相对较为可靠，但仍出现大批逃跑，“第一期地方上调去训练的差不多逃走了十分之六”。\footnote{《中共粤赣省委代表团通知第一号》，《闽粤赣革命历史文件汇集（1932～1933）》，第218页。}随着第五次反“围剿”失利，苏维埃政权处境不断恶化后，干部反水事例更多。1934年1月，“武平梁山乡主席于正月间叛变，并带去梁山游击队十三名，步枪十三支”。\footnote{《警钟——向着军区分区敲》，《红星》第26期，1934年1月28日。}3月，于都禾丰区“二百余群众，在少共区委组织科长的领导下，带着鸟枪梭镖等武器逃跑”。\footnote{《禾丰区破获反革命暗杀团》，《红色中华》第176期，1934年4月17日。}8月，“洛口区群众二百余在该区反动的区宣传部长、粮食部长、乡主席等的欺骗压迫之下向头陂逃亡”。\footnote{《陈伯钧日记（1933～1937）》，第267～268页。}


生存是民众的第一要务。由于国民党的封锁及前线供应的需要，后方生活日益紧张，生存受到严重威胁，民众不得不自寻生路，这也是出现逃跑的原因之一：

\begin{quote}

四十岁以上的男人很多都陆续地跑出苏区，到国民党区投亲靠友。有时搞到一点什么东西，也偷着回来一两次接济家里。因为他在家里实在是难以生活下去。农业上那些地方都是山地，种植业不发达，有的连种子都没有，又缺少食盐，基本的生活都没有办法保证。而我们也没有办法来解决这些问题。这种逃跑现象各县都有，特别是那些偏僻的山区里面，跑起来人不知鬼不觉。\footnote{《李一氓回忆录》，人民出版社，2001，第156页。}
\end{quote}


当时苏区之外也有人观察到：“彼方物质缺乏，盖亦足为离散群众之一原因也。”\footnote{陈赓雅：《赣皖湘鄂视察记》，第6页。}


1933～1934年苏区出现的群众集中逃跑现象，和当时中央领导错误指导有密切关系。查田运动、肃反、发行公债、借谷及扩红运动中的一系列问题酿成了苏维埃政权与群众间的紧张关系。当时中共中央领导人缺乏在政权、集体、个人间寻求利益平衡点的意识，一味强调无条件服从，夸大思想斗争的作用。时在湘赣苏区的萧克深有体会地谈道，1934年7月，湘赣苏区面临国民党军压迫，被迫向井冈山地区寻求退路，结果却很不理想：

\begin{quote}

由于“左”倾路线的错误，曾在这里实行过过左的社会政策，如地主不分田，富农分坏田，损害富裕中农利益，对反水群众不注意争取，对知识分子也以其成份作去留使用的标准等等。这样，就增加了我们工作中的困难……我们当时想恢复井冈山，可是我们上去后，连留在山上老百姓大部分都躲起来了。\footnote{萧克：《回忆湘赣红军》，《湘赣革命根据地》（下），第995～996页。}
\end{quote}


就在萧克说到的同一个月，在赤水、大寨脑一带作战的中央苏区红军发现，这里不少群众“加入刀匪，在板桥、在良田用石头打红军杀哨兵，摸哨等……因此战区粮食较为困难，而要由三十里许的驿前及六十里许的小松运来供给”。\footnote{《十五师大寨脑战斗详报（1934年7月24日）》，《中央苏区第五次反“围剿”》（下），第157页。}由于上述问题的严重性，以致毛泽东在二苏大报告中特别强调：“关于赤白对立问题，群众逃跑问题，食盐封锁问题，被难群众问题等，必须根据阶级的与群众路线，很好的给予解决。必须把造成赤白对立与群众逃跑的原因除掉了去。”\footnote{毛泽东：《中华苏维埃共和国中央执行委员会与人民委员会对第二次全国苏维埃代表大会的报告》，《苏维埃中国》，第302页。}鱼水关系变成了猫鼠关系，多年后读至此也不能不让人痛心。


值得特别指出的是，虽然当年年轻的中共领导人在艰困局面下，没有表现出驾驭危局的足够能力，导致社会政治危机丛生，但他们并没有试图掩盖问题，也一直在寻求解决的办法，正由于此，作为后人，我们还能从中共留下的文件中体会到当年危机的严重。这种负责任的坦率态度，值得后人予以尊敬。

\section{扩红与开小差}


和群众的逃跑相应，第五次反“围剿”期间，开小差成为困扰红军的重大问题。这一状况的出现，本身应为苏区一系列问题反映在红军中的结果，其对红军造成的伤害不可谓不大。


战争环境下，虽然红军严明的纪律与组织保证了红军旺盛的战斗力和团结精神，但开小差的现象在各个阶段都不可能完全消灭。初期，随着革命形势的起伏，队伍中开小差的比率起落不定。苏区巩固后，严重的开小差问题则主要由大规模扩大红军引起。早在1930年底，毛泽东在吉水东塘调查时就注意到：“本乡先后共去了七十九个人当红军，都是鼓动去的。但最后一批四十六人中，有四五个人哭着不愿去，是勉强去的。”\footnote{毛泽东：《东塘等处调查》，《毛泽东农村调查文集》，第259页。}扩红已经不能保证完全自愿。而且，“好些新兵到了分派到各师团去施队时，他们才知道是要他们当红军”。\footnote{《红军第一方面军前敌委员会江西省行动委员会通告（1930年11月11日）》，《中央苏区第一次反“围剿”》（《江西党史资料》第17辑），中共江西省委党史资料征集委员会、中共江西省委党史研究室编印，1990，第40页。}由于参军本非完全自愿，开小差也就很难避免。1931年底，红三军团第三师发现：“特务一连长、师部副官参谋、传令排长是党员，开小差走了。”\footnote{国平：《转变中三师的党》，《武库》第7期，1931年12月31日。}不过，这一时期，无论是扩红中的强迫还是开小差，都还只是个别现象，尚处可控状态。


随着时间的推移，中央苏区兵员需求越来越大，扩红的困难也在加重，邓颖超写下其参加扩红会议的亲身经历：

\begin{quote}

征求党员自动当红军，首先由参加人（中局及县委）发言解释鼓动，继由出席党员大会红军学校党员发言鼓动与欢迎。第一次鼓动发言后，回答是静默沉闷，继之二次鼓励，依然是静默沉闷，三次四次，经过半小时的鼓动工作，终无一人来报名当红军！最后即提出不当红军的原因的问题来讨论，很久很久，才在百卅余人中，涌出一句“因为没有执行优待红军十八条”一句回音来。此外就再无他语了。\footnote{邓颖超：《瑞金城区党员大会经过与教训》，《党的建设》第4期，1932年9月10日。}
\end{quote}


大规模扩红虽然有反“围剿”战争客观需要的成分，但当时也有人指出，由于苏区人力有限，合理调配人力十分重要。1933年初，福建省委代理书记罗明曾提出：“机械的规定……要扩大多少红军，这是不好的”，强调：“地方武装顽强的斗争，对整个发展的配合与对主力红军行动配合的力量，胜过出三百红军送到前方的主力。”\footnote{罗明：《对工作的几点意见》，《中央革命根据地史料选编》（中），第380～381页。}这应该说是针对苏区实际并集数次反“围剿”战争胜利的经验之谈。


苏区民众具有浓厚的宗族、乡土意识，参加地方武装、在家乡保卫自己，许多人十分积极，但离开家乡参加红军，却被不少人视为畏途。萧克回忆：“农民想的是打土豪分田地，即便拿了枪，也只愿意在本地活动，不愿远出，也不大愿当大红军，大概这就是所谓‘土共’（敌人报纸轻视地方党员和农军的贬词，我们有时也诙谐地借用）的特点。”\footnote{萧克：《朱毛红军中的农民军》，《朱毛红军侧记》，中共中央党校出版社，1993，第45页。}萧克的说法在当年的报告中得到证实：“一般农民乃至一般党员若叫他在本地当赤色队，打团匪土匪，他为着保护自己还很勇敢，一说到调他当红军，他就不愿意。”非常明显反映这种地方意识的是：“他们的逃跑关系红军的出发地，例如红军长汀打仗，永定籍的红军就逃跑。”\footnote{《刘伯坚关于闽西军校报告（1930年11月）》，《福建革命历史文件汇集》甲8册，中央档案馆、福建档案馆编印，1986，第185页；陈正、刘伯坚：《关于工作情况的报告（1930年11月7日）》，福建省档案馆藏91/2/404。}固然苏维埃政权应该尽力宣传打破地方观念，但一定时期内对现有事实的承认也是制定政策不能不考虑的因素。正如罗明指出的：

\begin{quote}

在敌人进攻比较紧张当中，应先抓到这个时机来扩大独立师、独立团和其它地方武装，先解决目前紧急的打击敌人进攻的问题，从这样的地方的紧急动员中来提高地方武装和群众的斗争情绪，从中来扩大红军……这是适应群众的斗争的情绪和目前的斗争的需要，不是什么对地方主义的投降。\footnote{罗明：《对工作的几点意见》，《中央革命根据地史料选编》（中），第379页。}
\end{quote}


不片面地追求主力红军的扩大，而是通过多种方式因地制宜地尽可能地增加苏区的地方武装力量，这是当时苏区人力、物力资源十分有限背景下相对有利的选择。当时报道反映：福建宁化“曹坊群众自动要求加入游击队打里田、四堡的反动派；中沙群众要求加入游击队打水西、安远司反动派，巩固中沙政权。而区委的同志不能抓紧群众的要求来解决这个问题，再进行各方面的详细的解释工作，只是很死板的‘到模范团去’，所以中沙的动员用欺骗的办法，结果到了宁化城完全开了小差，反而阻碍了扩大红军”。\footnote{《八月份宁化党扩大红军的转变》，《斗争》第29期，1933年10月7日。}


在看到罗明的地方性视野有其成立基础的同时，还应指出，扩大主力红军终究是红军的必由之路。如果没有主力红军的建设，仅仅愿在家门口作战的武装难逃乌合之众之讥。因此，中共中央对罗明建议的拒绝乃至批判虽失之褊狭，但不能认为罗明的方案就可以使一切迎刃而解。问题的核心不在其他，关键还是巨大的兵员需求和紧缺的人力供给之间的落差，而这又是由战争的严酷性所决定的。1933年8月至1934年7月中旬，苏区突击扩大红军达到11.2万多人。\footnote{《一年来扩大红军的统计》，《红星报》第54期，1934年7月22日。}超常规模的扩大红军，伴随的是普遍的强迫，苏区各地都存在“‘开大会选举’、‘乡苏下条子’、‘以没收土地恐吓’、‘指派’、‘拈阄’、‘受训练或守犯人’、‘前方不好可以请假’、‘以电筒干粮袋□子’……等方法来强迫、命令、欺骗、引诱、收买群众去扩大红军”。\footnote{《中共闽粤赣省委接受中央局扩大红军决议的决议》，《闽粤赣苏区革命历史档案汇编》第1辑，第230页。}其结果带来了十分严重的开小差问题。1933年底，仅瑞金一县逃兵就达到2500人，经过强制突击并枪决、捕捉部分屡次逃跑者后仍有八九百人。\footnote{《胜利的瑞金突击月》，《斗争》第43期，1934年1月19日。}据红军总政治部1933年11月的统计，“一军团补充区域到十一月十五日集中到区的是一六六三人，到补充师的只有七二八人，路上跑了九三五人。这儿还没有计算在乡村报了名根本未集中以及从乡到区逃跑的人数。大概算起来我们只集中了报名人数中的十分之三、四”。\footnote{王稼蔷：《为扩大红军二万五千人而斗争》，《斗争》第37期，1933年12月5日。}新战士到部队后的情况也不乐观：“一军团补充师在十二月份二十天内，共逃亡一三八名，占新战士百分之十，洗刷的老弱与阶级异己分子十二名。三军团补充师从八月到十二月份上半月共逃亡八七九名，占新战士百分之十五，洗刷的老弱残废六九七名，阶级异己分子九名，占百分之十二。五军团补充师从九月到十一月份共逃亡七八八名，占新战士百分之廿五。”\footnote{《严重的问题摆在补充师团的面前》，《红星报》第21期，1933年12月23日。}


开小差现象出现于从征集兵员到部队服役的各个阶段。一是在报名和集中过程中就有大批开小差的：“长汀模范团因为是被欺骗加入工人师，到瑞金集中时只剩三分之一，三分之二开小差走了。”公略动员2400多人参军，结果中途开小差的就有千人。“宁化模范团成千人送博生沿途开小差只剩二百余人。”\footnote{汉年：《工人师少共国际师的动员总结与今后四个月的动员计划》，《斗争》第24期，1933年8月29日；《中共公略县委七、八、九三个月工作总报告》，《江西革命历史文件汇集（1932年）》（二），第238页；《中共福建省委工作报告大纲》，《中央革命根据地史料选编》（上），第505页。}1934年5～7月，第一补充师逃跑317人，第二补充师逃跑44人，第三补充师逃跑265人。\footnote{《消灭逃跑现象来纪念“八一”》，《红色中华》第217期，1934年7月21日。}二是集中到部队后，仍然有相当多的人开小差。1934年9月，红五军团第十三师因逃亡、生病减员达1800余人，几乎占到该部总人数的1/3。红一军团在1934年7月根据不完全的材料估计，逃跑的也“达二百余人”。\footnote{《陈伯钧日记（1933～1937）》，第136页；聂云臻：《与逃亡现象斗争的一年》，《红星》第56期，1934年8月1日。}三是反复开小差情况严重。福建上杭才溪区的共青团员王佳，“在红军中逃跑了五次，也做了五次的新战士”。瑞金黄安区有一个人先后开小差六次，后又被征召，且被任命为排长，“几乎把一排新战士都断送在他手上”。\footnote{《足智多谋的智多星》，《青年实话》第2卷第3号，1933年1月29日；汉年：《工人师少共国际师的动员总结与今后四个月的动员计划》，《斗争》第24期，1933年8月29日。}四是军官和老兵也加入逃跑行列：

\begin{quote}

江西全省动员到前方配合红军作战的赤卫军模范营、模范少队在几天内开小差已达全数的四分之三，剩下的不过四分之一，所逃跑的不仅是队员，尤其是主要的领导干部也同样逃跑，如胜利、博生之送去一团十二个连，而逃跑了十一个团营连长，带去少队拐公家伙食逃跑。永丰的营长政委也逃跑了，兴国的连长跑了几个，特别是那些司务长拐带公家的伙食大批的逃跑。\footnote{《江西省苏维埃政府、江西军区总指挥部联合通令——关于模范赤少队开小差问题（1933年4月16日）》，《江西革命历史文件汇集（1933～1934年）》，中央档案馆、江西省档案馆编印，1992，第107页。}
\end{quote}


正规红军出现开小差同时，地方游击队逃跑现象也相当严重。1932年下半年江西出现多起地方武装逃跑反水事件：“于都整个一排人的叛变当靖匪；宜黄独立团一连政委领导三个士兵叛变投白军；乐安独立团三十一个士兵的反水到勇敢队里去；会昌独立团……有一连整个在连长政委领导之下，企图投到广东去当土匪；宁都城市的赤卫军里面大批流氓靖卫团的重要分子，混在里面被反革命领导暴动，包围县苏，布置罢市，缴东门外少先队的枪。”\footnote{《为加强和巩固地方武装发展游击战争的决议》，《中央革命根据地史料选编》（中），第644页。}“宜乐黄陂新游击队排长政委于十二月廿一日带所属人枪廿余，藉打土豪为名，降了崇三都敌人，以后接连又是东陂新游击队一百余人在八都去投降敌人。”\footnote{《警钟——向着军区分区敲》，《红星》第26期，1934年1月28日。}由于开小差者人数众多，对其开展工作也相当困难，龙岩、瑞金等地“都有开小差的同志联合起来与政府对抗”。洛口“地主武装及小差团”甚至合作攻进洛口街市，而“党政机关人员逃避一空，毫无应付办法”。\footnote{邓发：《开展反对开小差的群众运动》，《斗争》第14期，1933年6月5日；《陈伯钧日记（1933～1937）》，第272～273页。}


开小差问题的日益严重，除各地为完成扩红指标进行强迫命令，造成民众的抵触情绪外，各种各样的利益冲突也是导致这一现象的重要原因。从地方来看，相当部分地区对开小差现象没有予以充分重视，由于人力紧缺，各地常常在工作中把开小差人员作为有经验的骨干加以重用。江西、福建都有地方政府“替红军请假、寄路票给红军回家”，甚至“拿排长、队长官衔作引诱”，逃兵“可以担任当地重要工作”。\footnote{《中共闽粤赣省委关于消灭团匪与土匪问题给各级党部的指示信》，《闽粤赣革命历史文件汇集（1932～1933）》，第201页；《江西省工农兵苏维埃第一次代表大会对扩大红军的决议》，《中央革命根据地史料选编》（中），第607页；《中共闽粤赣省委关于消灭团匪与土匪问题给各级党部的指示信》，《闽粤赣革命历史文件汇集（1932～1933）》，第201页。}地方游击队为提高自身的军事水准也有意留用部队中开小差者，这些人员在“福建三分区游击队中占十分之七”。\footnote{《关于游击队工作，总政治部1934年1月5日训令》，《斗争》第42期，1934年1月20日。}同时，有些地方政权负责人“自己的老弟、儿子、侄儿等逃跑躲避而包庇他们，因此其它的队员也都跟着逃跑躲避”。\footnote{《瑞金党的道路，是全国苏区党的道路》，《斗争》第73期，1934年9月30日。}更有个别地区干脆把逃兵作为生财之道：宁都湛田区“政府卖路票给红军开小差”。\footnote{严仲：《三个月扩大红军的教训》，《中央革命根据地史料选编》（中），第665页。}正如王稼祥激烈批评的，某些地方干部“没有采取必要的办法对付那些组织逃跑的领导分子，与屡次开小差的‘专家’，相反的甚至有少数分子还在地方机关中工作，在军事部工作，在赤卫队当干部……大大的阻碍了红军的扩大与巩固”。\footnote{王稼蔷：《紧急动员——为扩大红军二万五千人而斗争》，《斗争》第37期，1933年12月5日。}


从红军本身看，随着队伍的迅速扩展，部队训练和政治工作常常滞后，相当程度上影响到部队建设。由于文化基础差，政治训练的作用难以充分发挥，部队政治程度仍然较低：

\begin{quote}

在三师全师党的代表大会中，举行了一次政治测验，受测验的人都是党内活动分子，测验的结果是：有把罗章龙认为现在中国革命重要领袖的，有把苏联和英国当为殖民地的，有把反帝大同盟当为反革命组织的，有主张反对游击战争的，有不知道日本帝国主义侵占东三省的。最惊人的有好几个连政委都把国际联盟当作全世界革命的参谋部。\footnote{《巡视三师的零零碎碎》，《武库》第7期，1931年12月31日。}
\end{quote}


同时，部队也存在一些不良现象，包括“不到操不上政治课不识字不做墙报宿娼争风打架上酒馆”，“各部队都有许多嫖娼的，赌钱的，抽鸦片烟的”。\footnote{国平：《转变中的三师的党》、《反对腐化》，《武库》第7期，1931年12月31日；第9期，1932年1月21日。}一些地方游击队质量尤其难以保证。宁化县独立游击支队营“内容腐败，吃大烟、嫖姑娘的很多，纪律松弛”。\footnote{《霍步青给稼蔷同志信（1932年11月8日）》，《闽粤赣革命历史文件汇集（1932～1933）》，第273页。}福建莆田游击队更为荒唐，“因没有钱吃饭，化装反动军队去抢商店，结果给农民缴去枪六枝，被捕队员七人”。\footnote{《中共福州中心市委书记陈之枢给中央的报告（1933年6月19日）》，《福建革命历史文件汇集》甲13册，第91页。}虽然，总体上看这只是个别现象，但仍反映出一支纪律严明部队的建设需要持之以恒，不可能一蹴而就。


讨论开小差问题，还有一点不能不说到的是红军生活的艰苦。根据地初创时期，红军通过打土豪可以筹集大笔款项，随着打土豪财源的枯竭、根据地经济上的困窘，红军经费日益紧张，供应短缺的问题日益严重：“红军是完全没有储粮，单靠杂谷吃饭的士兵，每月只有三元的伙食费。在今天米（每元七升）与食品价高的时候是很苦的，大多数不独衣装成问题，并且半数以上的士兵连草鞋都没有而赤足作战的……士兵愿当赤卫队而不愿当红军。”\footnote{《家、弼、霖自闽粤赣苏区来信（1931年3月23日）》，《闽粤赣革命历史文件汇集（1930～1931）》，第72页。}中共报告中也痛心地承认：“士兵生活苦到万分。”\footnote{《中共湘赣苏区省委报告》，《湘赣革命根据地》（上），第599页。}在战争背景下，加入红军随时面临着牺牲的危险，又要承受艰苦的训练和生活，不少人望而生畏，闽西报告：“一般群众到红军去都为了生活艰苦而不去。”\footnote{《闽西同志口头报告（1930年12月19日）》，《福建革命历史文件汇集》甲16册，第110页。}当年的档案文件清楚记下了民众的所思所想：

\begin{quote}

过去第四军回到闽西，只永定一县就不少农民自动加入红军，后四军到江西去，经过几个月都请假回籍了，他们在群众中传达红军的生活如何艰难，我们去动员群众加入红军，他们首先就是障碍，无论怎样解释他们都不愿意再当红军，反转说，你们没有当过红军，哪里知道，我自己当过红军此系知道得清楚……每次征调红军，出于自动加入的，可以说是没有的。\footnote{陈正、刘伯坚：《关于工作情况的报告（1930年11月7日）》，福建省档案馆藏91/2/404。}
\end{quote}


艰苦的生活，使部分加入红军者难以承受，有逃兵承认：“红军每天要走一、二百里路，我们实在拖不得。”\footnote{《远关于湘鄂西苏区敌军和我军情况及红军中政治工作缺点给中央的报告》，《湘鄂西苏区革命历史文件汇集（1931～1934）》，第395页。}王稼祥等也意识到：“红军生活这样艰苦，与士兵的逃跑自然是有关系的。”\footnote{《家、弼、霖自闽粤赣苏区来信（1931年3月23日）》，《闽粤赣革命历史文件汇集（1930～1931）》，第72页。}


应该承认，在苏区的现实环境下，相当部分新加入的红军并没有非常清楚的政治观念，老红军李林写于1940年代的一份思想汇报作为原生态的材料，颇值一读：

\begin{quote}

本人出身农民，小时候在（也）受苦。家中无钱读书，十二岁在家耕田，帮人放牛。到十六岁至十八岁，这时候的思想就不想在家受苦，帮人放牛心里都想悲脱（背着）出来当兵，后来我母亲设法我跟我叔叔（继父）做小生意。当时我就不愿意跟我叔叔在一块，因为他经常打骂我，我不愿跟我叔叔做事，经常想出来当兵，面（免）以和我叔叔吵骂。以后红军来到我家，就参加本地游击队有半年。后来被地方反动捣乱了，那游击队分散了。我有（又）找正式部队，来时我母亲不让我出来当兵，以后队伍开在我家过时，我就跟队伍走，我母亲就躺在路边哭。当时看到那样哭，有（又）回去了一次，安慰我母亲不要着（生）气，我到外面不会吃亏，有人会帮助，叫她在家放心。同我母亲说这些话，我母亲都不让我出来，后来我就不管家的事情好坏。这就是我出来的思想情形。

……

到部队后的思想没有什么表现，后来就加入共产青年团员。加入共产青年团后，人家告诉我说，加入青年团为无产阶级谋利益，到共产主义社会。以后就有人问我，为什么要加入共产青年团？我就照上面跟照（着）说：要为无产阶级谋（利益），打土豪分田地。开始对于这些了解还不清楚。后来在党内受到教育才懂得些。

一九三四年三月加入青年团，三六年在定边转党……一九三四年在一方面军一军团电台当护员，那时候的思想有吃就算了，其他不管。\footnote{李林：《历史思想自传》，《建宁党史资料》第3辑，中共建宁县委党史办公室编印，1989，第119～120页。原文因错别字较多，引录时已在括号内做过修订。}
\end{quote}


读到这一段文字，一个老红军质朴的形象可以说呼之欲出，其实，这样的思想经历为很多人所拥有。从马克思主义经典作家所描述的无数个成长中的“马铃薯”组成红军这一事实，就可以理解艰苦环境下出现开小差等问题的不可避免性。


另外，从地域因素看，赣南、闽西偏处一隅，历史上经历战乱较少，民众从军热情不像长期经历战争地区那样高。籍贯湖南的邓文仪回忆，他出外当兵时，“大多数亲友都很称赞，并且轮流请饭饯行”，\footnote{邓文仪：《冒险犯难记》（上），台北，学生书局，1973，第21页。}这种现象在赣南、闽西很难见到。林彪曾将江西和湖南相比，谈道：湖南你再说不行，湖南人他愿意当兵。我们这个队伍到湖南就扩大了队伍。我们队伍一到江西，没有一个江西老百姓愿意当兵。我们到福建，福建当然有苏区，但是也没有一个人愿意当红军。\footnote{陈毅同志“九一三”以后的讲话，1971年10月下旬。林彪此言，当然说的是一个总体感觉，中央苏区也不是完全没有志愿参军者，如吴法宪回忆就提到：“我看到有的同乡，只比我先参军几个月，就当了副班长，我看红军当官很容易，我想当官。”见《吴法宪回忆录》（上），香港，北星出版社，2006，第22页。}


中央苏区发展过程中，尤其是第五次反“围剿”期间出现的红军、游击队出现的开小差现象，是苏区一系列社会政治难局下形成的困境，根源在于战争的残酷和苏区人力资源的缺乏。中共中央指导方针的失误固然加剧了这一现象，但只是导致问题出现的诸多因素之一，不能要求其承担起全部的责任。这一点，在赤白对立现象的形成发展过程中，也是如此。

\section{“赤白对立”}


苏维埃时期，苏区周边地区曾出现严重的“赤白对立”现象。所谓赤白对立，是指苏维埃区域与非苏维埃区域之间的对立，它不是由苏维埃革命加剧的阶级间的对立，而是一种非阶级的由多种因素引发的以地域为中心的冲突，主要发生在苏区边境地区。


早在1928年10月，赣东北就有报告提到：“环绕割据区的民众，还不知道我们的好处。土劣已感觉我们与他们不利，设法使民众起来反抗我们了。故环绕割据区域的民众非常反动，每日跟着反动军队，来我地抢东西。凡民众有食器用，只要能搬运者，莫不抢劫一空。”\footnote{《弋阳、横峰工作报告（1928年10月）》，《江西革命历史文件汇集（1927～1928年）》年，第296页。}1929年1月，滕代远向中共湖南省委报告，江西“铜鼓的民众因被平江游击队烧了很多房子……当红军到铜鼓县城时，所有男女老幼各种货物桶，一概搬运走了，不但找不到党的关系，连饭也没得吃”。\footnote{《滕代远向湖南省委的报告（1929年1月12日）》，《湘鄂赣革命根据地文献资料》第1辑，第47页。}这是我们可以看到的较早谈到这一问题的报告。


赤白对立的出现，和中共阶级革命的宗旨存在距离，作为苏维埃革命的倡导者，中共以领导全国人民革命为己任，非苏区区域从原则上说，理应是革命的发展对象，当地人民潜藏的革命热情和苏区人民也应是一致的。但是，赤白对立现象却在相当程度上挑战着中共这一理念，严重影响着苏维埃区域的发展：“各县警卫营或连，特别是赤少队很多都不愿到白区域去，以为白区群众都是些‘虎豹豺狼’，不能同它们接近……有许多听到要他去白区工作他就哭起来，甚至哭得饭也不想吃。”与此同时，“白区群众的怕游击队名之为‘刀子队’，造成了赤白对立的现象，如铜城铁壁一般”。\footnote{《中共河西道委给苏区中央局的综合工作报告（1932年5月17日）》，《江西革命历史文件汇集（1932年）》（一），中央档案馆、江西省档案馆编印，1992，第111页。}


这样一个中共不希望看到的现象，其最初出现，客观看，和中共苏维埃革命初期实行的错误政策不无关系。当时有关文件详细分析了造成赤白对立的原因：“因我们的工作不好，特别是盲动主义引起白区群众的反感”；“旧社会遗留下来的地方观念，宗族观念与发生过械斗”；“出于豪绅地主反动派的欺骗造谣，挑拨离间”。\footnote{《湘赣省军区政治工作会议决议案》，《湘赣革命根据地》（上），第685～686页。}这里，首要的就是盲动主义影响。1928年前后，中共在盲动政策指导下，普遍执行了烧杀政策，对地主等革命对象进行肉体消灭，打击对象甚至扩展到苏区外的一般群众，中革军委总政治部指出：

\begin{quote}

红白两边，杀过来，杀过去，成了不解的冤仇。这其中，在革命方面犯了许多盲动主义、报复主义的错误，乱抢乱烧乱杀的结果，反造成那边的群众更加坚决的反对革命。有一次工农群众和赤卫队打到龙聚区的吴公山，一烧就烧一千多家房子，这样越发使他们接受豪绅地主的欺骗，反对革命了。\footnote{中革军委总政治部：《争取三都七保的意义和工作方法》，《中革军委总政治部通讯》第3期，1931年2月26日。}
\end{quote}


同时苏维埃革命政策本身也必然严重触及地主等农村有力阶层的利益，作为报复，受到国民党军支持的反苏维埃地方武装回到当地后，往往对参加革命的民众进行屠杀。湖南平江自1928年初暴动后，“杀戮豪劣和反动份子，计在数千，而同志和革命民众殉难的，亦不下数千人”，许多地区“数十里或百数十里，几无一栋完善的房屋，无一处尽青的山，共计全县被烧的房子，总在十分之四、五”。\footnote{《滕代远向湖南省委的报告（1929年1月12日）》，《湘鄂赣革命根据地文献资料》第1辑，第43～45页。}江西上犹在红军退走后，土豪地主组织的民团对民众“不问首犯盲从，一律处以死刑，其中遭冤枉而死者，不知凡几”。\footnote{《上犹人之赤祸谈》，《共匪祸赣实录》第2期，第58页。}地主疯狂的屠杀又激起民众的强烈愤恨，以致“报复心理非常浓厚，盲动主义时代精神的复活，群众无论如何要求以烧杀抢劫来答复白色恐怖，其气之高真不可制止”。\footnote{《赣西南刘作抚同志报告》，《中央革命根据地史料选编》（上），第225页。}中共有关文件明确谈道：“各苏维埃区域边境严重的红白对立现象，就是这种报复主义造成的结果。”\footnote{《中共苏区中央局通告第十号，地方武装的策略组织和工作路线》，石叟档案008·5524/2612/0559，中国社会科学院近代史研究所藏。}


对立情绪不断蔓延，相互间的报复行为，常常超越阶级对抗的范畴，变成区域之间的对抗。福建“蛟洋农民烧丘坊房子二百余家，白砂赤卫队烧茶地房子九十余家”。宁德横坑民团与中共开辟的游击区敌对，引起苏区群众愤慨，“见横坑人即杀，横坑人就不敢向游击区域来买东西”。\footnote{《中共闽西第一次代表大会之政治决议案》，《中央革命根据地史料选编》（中），第117页；《中共福安中心县委工作报告（1933年7月10日）》，《福建革命历史文件汇集》甲19册，第120页。}同时，这一状况的出现，和中共开展革命时为鼓动民众、克服地方观念，常常组织农民跨村跨乡打土豪有关。跨村跨乡活动在苏维埃区域内虽有可能触发宗族间的冲突，但总体处于可控状态，且村庄间的运动是相互的，负作用不明显；在边区则变成苏区对白区的单向运动，且由于对当地具体情况不了解，打土豪行动往往失控，有些部队打土豪“离开阶级的（标准），以有猪有鸡有谷有鱼为标准的乱打，结果把有猪有鸡有谷的中农和贫农也打了，破坏与中农的团结和阶级利益”。\footnote{《彻底执行明确阶级路线与充分群众工作》，《努力》第5期，1933年8月6日。}有些“地方党部及政府不仅不纠正此种错误，而且造成理论说这是赤白区域对立，他们不了解赤白区对立，是赤白区两个政权的对立，而不是赤白区两地群众对立，他们把地域的对立，代替了阶级的对立”。\footnote{《邓中夏同志关于红二六军的报告》，《湘鄂西苏区革命历史文件汇集》第2册，第36页。}这种不顾阶级关系的盲目报复和掠夺行为，虽然和革命对象一方的白色恐怖及压迫有着重要关系，却客观促成了赤白对立。


作为一个以群众革命为生存基点的政党，中共对群众利益、群众情绪始终予以高度重视。赤白对立现象出现伊始，中共各级部门已有所意识，随着其范围的扩大和危害的增长，更予以高度重视并力图加以克服。1931年第二次反“围剿”前后，红一方面军总前委严厉批评了红四军“在水南烧土豪房子时烧了一条街”\footnote{《总前委第六次会记录（1931年6月2日于建宁）》。}的错误，责令公开作出赔偿。当计划进攻和赤区有强烈对立倾向的七坊时，更明确要求：“到七坊后要开和平会，立和平公约，第一条要两方不再相打，大家一起打土豪。”\footnote{《总委第九次会议记录（1931年6月22日于康都）》。}中革军委总政治部也指出：

\begin{quote}

坚决反对烧杀政策和报复主义。游击队在每次进攻反动统治区域之前，必须对自己部队及参加斗争的群众宣传纪律，绝对禁止侵犯贫苦群众的利益，不准乱烧他们的屋，不准乱杀他们的人，不准乱拿他们一点东西……豪绅地主的谷子衣服猪牛用具，原则上要完全发给当地贫苦群众（号召当地群众去夺取），这是发动群众的斗争的必要策略。游击队和参战的红色群众决不可取得太多。\footnote{《中共苏区中央局通告第十号，地方武装的策略组织和工作路线》，石叟档案008·5524/2612/0559，中国社会科学院近代史研究所藏。}
\end{quote}


中共采取的这些措施在一定时间和程度上收到了成效，但是一直到苏维埃革命终止，赤白对立现象也未得到完全制止，在一些地区、一些时段甚至有愈演愈烈之势。之所以如此，一方面是因为早期赤白对立的后患难以短时期消除，其负面影响会持续存在；另一方面，苏区内外造成赤白对立的新因素不断出现，使赤白对立成为大部分苏区难以克服的硬伤。


一部分地方武装和游击队乱打土豪是赤白对立持续并加深的重要原因。苏区大量的地方武装供应依赖地方，在战争环境下常常难以为继，利用到白区活动机会，筹集粮款，是其解决自身生存问题的一个重要途径，而各级机关对这一做法实际也采取默认态度：各级部门“检查工作时，首先问‘多少钱’，不去检查打豪劣地主与民众的情绪如何”。\footnote{《中共福州中心市委符镭关于福州工作情况给中央的报告（1933年12月）》，《福建革命历史文件汇集》甲13册，第196页。}有的游击队把打土豪收入“抽十分之一作伙食尾子分”。\footnote{《中共南广县委给苏区中央局的综合报告（1932年10月）》，《江西革命历史文件汇集（1932年）》（二），第98页。}同时，对于游击队到白区活动，中共中央指出：“分发豪绅谷物衣服给群众，是挺进游击队在数小时内取得群众拥护的最有效手段，每到一地，必须在最近期间最大限度进行这一工作。”\footnote{《江西军区政治部关于坚决执行中革军委“游击队怎样动作大纲”中挺进游击队之工作任务》，《江西革命历史文件汇集（1933～1934年）》，第364页。}这就是说，游击队进入白区后，必须迅速展开打土豪，在对情况不熟悉背景下，由此造成的一些问题具有相当的必然性：“地方武装过去行动，大多数都是陷于单纯筹款的泥坑中，许多行动不是为了群众利益而是自己去找经费，一到白区，豪绅地主走了，贫苦工农也乱捉乱打一顿，造成一种白区群众害怕游击队，甚至在豪绅地主欺骗之下来反对苏区，为难游击队，造成一种脱离群众的赤白对立的严重现象。”\footnote{《中共湘赣苏区省委报告（1933年2月）》，《湘赣革命根据地》（上），第594页。}后来，毛泽东在回顾这一段历史时，曾痛切谈道：“土地革命时期打土豪办法，所得不多，名誉又坏。”\footnote{《毛泽东年谱（1893～1949）》（下），第213页。}


除维持生存的因素外，地方武装和游击队军纪、政治工作薄弱、素质较差也是导致上述问题的重要原因。地方武装发展过程中“总带有或多或少的强迫，甚至完全用强迫命令而编的”，“兼之平时没有教育训练与党的领导薄弱……到白区乱抢东西则是普遍的发现”。“因为没有明确的阶级路线，有许多赤色武装到白区去，不坚决实行三大纪律八项注意，甚至负责人也不切实做到，工农的东西乱拿，群众跑出去了，以为这些是‘反革命的东西’，也可以打土豪！”\footnote{《中共会昌县委十、十一两月工作报告（1932年12月3日）》，《江西革命历史文件汇集（1932年）》（一），第367页；《中国工农红军湘赣省军区总指挥部白区工作大纲（1932年6月）》，《湘赣革命根据地》（上），第319页。}以致指导机关明确要求“党要坚决转变乱拿白区许多东西，不问穷人富人都捉来的许多错误，应当把赤卫军、游击队在白区行动的政治教育与军事纪律的建立成为目前主要工作，才能避免赤白对立”。\footnote{《中共鄂豫皖中央分局给陂孝北县委会的指示信》，《鄂豫皖革命根据地》第1册，第397页。}


随着苏维埃区域的扩大，当其向城镇扩展时，传统的城乡间相互敌视也对赤白对立现象发生影响：“城市方面的一般人，对于乡下人有鄙视欺侮的举动，是很普遍的事实，因此发生了城乡恶感。”\footnote{《闽西斗争意义与教训的讨论》，《福建革命历史文件汇集》甲4册，第27页。}“农民上县的时候，则检查异常严厉，有时借故寻衅，将农民捉去罚款，弄得一般农民不敢到城里去购买货物，农民非常怨恨城市。”\footnote{《鄂东巡视员曹大骏的报告（1929年8月31日）》，《湘鄂赣革命根据地文献资料》第1辑，第125页。}由于此，“在游击战争发展到城市去，农民便摧毁城市以泄恨，少不免影响中小商业停闭”。\footnote{《苏维埃区域的经济问题》（《红旗》编者按），《中国苏维埃》，三民公司，1930，第53页。}闽西中共地方苏维埃组织农民攻城时，“农民更喊着‘杀尽城内人’、‘烧尽城内屋’的口号。那时农民痛恨城内人的心理是十分一致的。他们都说：打进了城不但要杀尽抢尽烧尽，而且还要将城墙拆去。于是城内豪绅地主便利用这些口号去煽动城内贫民仇视我们。果然，一般贫民受其利用，做侦探、当团丁、做向导，无所不为”。\footnote{《闽西斗争意义与教训的讨论》，《福建革命历史文件汇集》甲4册，第27页。}


应该承认，在所有上述因素中，农民本身是造成赤白对立的基础性原因。中共的苏维埃革命和农民的支持息息相关，为获得农民的支持，一定程度上对农民利益的让步不可或缺。事实上，当中共武装攻打城市或到白区活动时，当地群众通常都是积极的参加者。这一方面反映了农民的革命要求，同时也和其搜罗财物这一利益目标有关。在此背景下，不分阶级、贫富乱拿财物变得难以避免。北上抗日先遣队报告：“我们打罗源县时，群众戴了小斗篷，拿了梭镖、扁担，足站了有五里路长。县城一打开就自动地拿东西，阻也阻不住，话也听不懂。”\footnote{《红军抗日先遣队北上经过的报告》，《闽浙皖赣边区史料》，中共中央宣传部党史资料室编印，1954，第52页。}福建方面也反映：“贫民的房子也去搜，上至银钱宝贵的物件，下至很破很败的衣服，不管是贫民、雇农、中农都把它掠回来。发动群众去参加就是去掠东西。有的群众说‘你们得的是钱，我们得的是坏东西’。攻进乡村时豪劣地主逃跑了，存下来小儿老人，也杀得鸡犬无存。”\footnote{《中共福州中心市委符镭关于福州工作情况给中央的报告（1933年12月）》，《福建革命历史文件汇集》甲13册，第196页。}江西万载农民甚至出现“七八间（处）抄抢队的组织”。\footnote{《湖南省委巡视员蒋长卿巡视湘鄂赣边境的报告（1929年12月20日）》，《湘鄂赣革命根据地文献资料》第1辑，第193页。}


当时，中共领导人对如何在发动群众参加斗争同时又保持严格的纪律颇感困惑，默认群众的抢掠行为势必破坏苏维埃政权的形象，强硬制止又担心引起群众不满。当中共刚刚开始军事斗争，南昌起义军退到汕头时就遇到了这样的问题，当时中共采取了严厉措施加以制止，但不久中共中央就改变了这一做法，并对执行制止方针“逮捕并杀乘机抢劫贫民三人”的责任人“以留党察看一年之处分”。\footnote{《政治纪律决议案（1927年11月14日）》，《中央通信》第13期，1927年11月30日。}然而放任不管造成的混乱其实也是中共一个沉重的政治包袱，因此，当1927年底的放任方针某种程度上和“左”倾盲动联系在一起后，中共也在调整自己的政策，1930年赣西南报告称：

\begin{quote}

（农民自卫军）五军攻下分宜后即准其进城，但进城后则全不问贫富先抢劫一空，有时还乱杀乱烧，五军因得了这一经验打下袁州即不许进城，同时还向他讲演并说明不许进城是因为敌人武装还没有完全缴得，恐进城受误伤，并马上没收一部分东西分给他们。但在当晚约一千余人爬进城来抢劫了数十家烧了房屋，军部马上派人来制止无效，继派武装弹压无效，并抢弹压士兵的枪，士兵以正当的防卫向空中开枪示威误打死一农民群众即镇压下去了，但引起了农民的反感。\footnote{《赣西南刘作抚同志报告》，《中央革命根据地史料选编》（上），第263页。}
\end{quote}


当地负责人对这一处理及其后果没有把握，请示中央“以后再逢有这样的事是如何处理”。\footnote{《赣西南刘作抚同志报告》，《中央革命根据地史料选编》（上），第263页。}确实，完全放任会导致秩序失控，但加以约束又会影响到农民积极性，中共在这一问题上颇有点左右为难。1931年，红一方面军总前委在计划进攻七坊时明确要求“群众赤卫队绝不要去”；\footnote{《总委第九次会议记录（1931年6月22日于康都）》。}苏区中央局也激烈批评：“各地破坏城市尤其是在城市中大烧房屋的办法，完全表现流氓路线，农民意识失败主义是非常错误的，以后要极力纠正。”\footnote{《中共苏区中央局通告第十号，地方武装的策略组织和工作路线》，石叟档案008·5524/2612/0559，中国社会科学院近代史研究所藏。}但总的来看，中共更多时候是采取教育和引导的方式，不加以过于激烈的控制，由于农民本身的利益冲动，抢掠行为仍无法完全制止。福建连江群众到白区去，拿走的东西“从棉被衣服直到饭碗，火钳都要被他们带着回去。而且还说：‘我跟你们打土豪，你们是很划算得来的，我们只拿一点东西，但是你们却由我们的帮助罚了很多钱和得了很多武装”。\footnote{《中共福州中心市委关于巡视连江工作报告（1933年11月10日）》，《福建革命历史文件汇集》甲3册，第148页。}


除中共本身原因外，中共的对手方国民党方面及地主也为赤白对立的形成推波助澜。毫无疑问，在动员农民的能力、方法、投入的精力及可利用的资源、手段上，国民党在苏维埃时期远远无法和中共相比。于国民党而言，赤白对立可以有效地限制中共力量向其控制区域的渗透，而其对苏区的影响，由于拿不出像土地革命这样富有号召力的实际措施，本来就困难重重。因此，和中共极力想消除赤白对立不同，赤白对立为国民党统治地区构筑了一道天然的屏障，是他们乐于看到的。同时，维持及造成赤白对立，也符合其封锁苏区、打击中共的战略目标。


国民党方面制造和加剧赤白对立主要依赖的是苏维埃革命的被打击者及中共在开展革命过程中的一些错误。随着苏维埃革命的进行，苏区内外的地主、富农作为革命的打击对象成为国民党政权的坚定拥护者，同时一部分中农及包括贫农在内的其他一些阶层出于对中共的误解也产生恐惧心理，而中共苏维埃革命初期的盲动政策及后来的“左”倾政策都加剧了这一倾向。当国民党方面试图在政治上与中共展开竞争时，这些都成为中共可被突破的软肋。和国民党政权一样，作为革命中的被打击对象，赤白对立也符合着豪绅地主的利益，因此，他们比国民党政权更积极地制造着赤白对立：“龙港的豪绅（非赤色区域）利用宗法社会关系，鼓动一些盲目群众，准备向当地同志进攻。”\footnote{《鄂东巡视员曹大骏的报告（1929年8月31日）》，《湘鄂赣革命根据地文献资料》第1辑，第131页。}


当国民党政权和豪绅地主在赤白对立这一问题上达成一致时，他们最常动用的资源是宗族关系。主要由地主构成的士绅阶层在农村社会具有重要的影响力，由于他们在乡村中拥有的财富、文化、社会资源，通常成为宗族的控制者和代言人。苏维埃革命展开后，为对抗中共革命的影响，国民党政权和地方豪绅充分利用宗族制度并将其与地方观念结合，发挥出相当的作用。福建漳州报告：“这里姓杨的农民，我们没有工作，在士绅地主领导之下帮助反军进攻我们，这是给我们火线上一个很大的打击。”\footnote{《北冀关于漳州红三团行动的报告（1933年8月25日）》，《福建革命历史文件汇集》甲11册，第91页。}国民党方面通过宗族和地方观念的号召，在一些宗族和地主豪绅力量较强地区，形成制造赤白对立对抗中共的有力力量。江西兴国、于都、宁都、永丰四县交界的三都七保地区民性“在历史上有名的蛮悍，从来不纳税，不完粮，不怕官兵”，苏区建立后，他们“受土豪劣绅的欺骗，中氏族主义的毒很深。那些豪绅地主团结本姓穷人的口号是‘宁可不要八字（命），不可不要一字（姓）’，这种口号在那些地方有很大的影响，因此所有的群众都被豪绅地主抓在手中”。\footnote{中革军委总政治部：《争取三都七保的意义和工作方法》，《中革军委总政治部通讯》第3期，1931年2月26日。}对此，何应钦曾报告，该地“民情最为强悍，反赤三年，赤匪受损甚大，视为赣南赤区最大障碍。卒因种种诡谋软化各区，仍不敢用高压手段，而赖村圩、汾坑圩、马鞍石等处至今尚守寨不屈。即被击破之村圩，多数民众仍持怀报复之心，是以我军一到，良民大半归来，热烈欢迎”。\footnote{《何前敌总司令官代电宁于兴永四县交界之三都七保人民反赤情形》，《陆海空军总司令部行营党政委员会会报》第5～8期合刊，1931年8月31日。}


在挑起赤白对立的过程中，农民好利的心理常被国民党所利用。如福建漳州民团“配合各地如潮水般的反动群众……向赤区进攻，抢掠赤区群众的猪、牛、粮食，搬不动的东西放火烧，锅子不要的就打破”。\footnote{《中共漳州县委书记蔡协民给厦门中心市委的工作报告（1932年8月16日）》，《福建革命历史文件汇集》甲19册，第24页。}1933年第五次“围剿”前夕出任国民党泰和县长的帅学富回忆，他在组织由苏区逃出的难民到苏区抢劫时公然声称：

\begin{quote}

你们挑选年富力强壮丁，手持梭标，身背匾担，跟在我保卫团后面前进，俟我打进匪区村落后，由你们抢劫，见牛牵牛，见谷挑谷，可是抢来任何物资，不准私藏己有，都由你们委员会，作公正合理的配给每一个难民享受，得来枪支，亦交你们义勇队使用……从此我这个县长，成为打家劫舍的强盗头子了。\footnote{帅学富：《五车书室见闻录》，台北，文海出版社，1981，第152页。}
\end{quote}


作为一种贯串苏区发展始终的现象，赤白对立的产生、延续，除了前文已经谈到的诸多现实原因外，更进一步看，它还和苏维埃革命的起源、动力，中国农村阶级分化的实际状况密切相关，这些因素和现实的原因交相影响，既成为一些导致赤白对立现象的政策得以出现的内在原因，又使得一些错误政策的负面影响被加深、放大，从而进一步加剧着赤白对立的发生、发展。


苏维埃革命是中共在国共合作破裂，自身面临生存危机时的选择，军事的推动是苏维埃区域形成、发展的主要支配力量，由武力所造成的苏区与非苏区的分隔，使区域的对立极易成为现实。苏维埃区域多在交通阻隔的山区，这里的特殊地理和经济状况影响着大地主的发育，以自耕农为主的农村社会格局，对中共阶级革命的判断和实践带来困惑，相当程度上成为中共过火政策的重要诱因；而这些地区特别明显的公田制度，和根深蒂固的宗族制度相结合，为赤白对立的发生提供了社会条件。农民追求利益的天性，决定了中共发动农民过程中既可以充分运用利益驱动号召农民，同时也可能需要承受这种利益冲动带来的消极影响乃至重大破坏。


赤白对立在诸多因素影响下，成为苏区的一种痼疾，它的存在，事实上成为苏区发展壮大的绊脚石。由于赤白对立，“经济不能流通，不但小商人不能做生意手工业的不能出售，而一般农民日用必需之品（油盐等），也因此而缺乏甚至买不到，因为豪绅地主及大富农都被赶出去了。同时又因抗债的关系，农民无处借贷，粮食也不能出售，所以在赤区农民感觉革命后更痛苦，虽然是没有地主豪绅及高利贷的压迫和剥削了”。\footnote{子修：《赣北工作综合报告（1930年7月）》，《江西革命历史文件汇集（1930年）》（一），第264页。}这一描述当然不一定完全准确反映苏区的实况，但确实说出了赤白对立对苏维埃政权巩固、发展的障碍。更重要的，赤白对立严重影响到苏区的对外发展，在赤白边境地区，由于“侵犯了贫苦工农的利益，以及豪绅地主武装的镇压与欺骗，我们部队一到该地，常常上山打埋伏，不同我们见面”。\footnote{《中共湘赣省委第二次代表大会工作报告》，《湘赣革命根据地》（上），第644页。}而游击队所作所为，使“白色区农民欢迎红军不爱游击队，甚至要求打游击队”。\footnote{《鄂豫边革命委员会报告（1930年4月）》，《鄂豫皖革命根据地》第2册，第89页。}在福建就有“群众要和土匪一起的拿长筒火炮来打我们”，\footnote{《中共闽粤赣省委关于领导和参加革命战争给各级党部的指示信》，《闽粤赣革命历史文件汇集（1932～1933）》，第187页。}安远、南丰县的白区边界的群众，“到处都向我们打枪，捉杀我们的红军病兵，及落伍士兵”。\footnote{《为加强和巩固地方武装发展游击战争的决议》，《中央革命根据地史料选编》（中），第645页。}由此在赤白边境地区形成一种“白打赤、赤打白依然部落式的战争”。\footnote{项英：《闽西的一般政治情形》，《福建革命历史文件汇集》甲8册，第256页。}


在赤白对立影响下，红军前出到边区乃至国民党区域作战，也面临着群众支持的问题。中央苏区第五次反“围剿”初期，红军在赤白边境地区作战时就遇到“军队打仗群众旁观，请不到向导，弄不到担架，脱离群众不发动群众积极性的怪现象”。\footnote{《彭滕关于我军今后作战的意见（1933年10月24日）》，《中央苏区第五次反“围剿”》（上），第94页。}第五次反“围剿”前夕彭德怀率部在闽西作战时深有感触：

\begin{quote}

群众对我们的态度也不热烈，召集群众大会，只有很少的人参加。这使我开始感觉到根据地内的土地政策有问题：地主不分田地，逃到白区流窜；富农分坏田，也有不少外逃；在加紧反对富农的口号下，打击了少数富裕中农，也有外逃者。他们伙同散布各种坏影响，使得边区工作很不好做。赤白对立，经济封锁，越来越严重。\footnote{《彭德怀自述》，第181页。}
\end{quote}


赤白对立使中共的阶级革命方针常常遭遇到某一地区民众多数的抵制而难以发挥作用，以致有人对白区民众丧失信心：“以为白区群众都是反革命的，与白区交通完全断绝，故意的形成赤白对立。”\footnote{《中国工农红军湘赣省军区总指挥部白区工作大纲（1932年6月）》，《湘赣革命根据地》（上），第318页。}


苏维埃革命期间，对阶级革命已经熟手的中共，却在赤白对立这样一个非阶级现象中显得有些应付乏策，这显示了历史进程的复杂，也指示着中共革命不仅仅应该正面应对阶级问题，还要客观面对社会现实，注意到另外一些更为复杂的社会政治现象。赤白对立的发生、延续并不简单是一个政策错误问题，而应有着更深刻的社会政治原因，物质资源和政治资源的纠结与选择，是中共在赤白对立问题上举棋难定的关键。还是那句老话，在高歌猛进的革命大纛后面，柴米油盐总是会顽强地显示着它们的存在，革命要从理想和浪漫中向前推进，依然离不开一点一滴的改造之功。

